\documentclass[11pt,a4paper]{article}

\usepackage[T1]{fontenc}
\usepackage[utf8]{inputenc}
\usepackage[italian]{babel}
\usepackage{geometry}
\usepackage{array}
\usepackage{booktabs}
\usepackage{longtable}
\usepackage{xcolor}
\usepackage{enumitem}
\usepackage{hyperref}
\usepackage{titlesec}

\geometry{margin=2.2cm}
\raggedbottom
\setlength{\emergencystretch}{3em}
\hypersetup{
  colorlinks=true,
  linkcolor=blue!55!black,
  urlcolor=blue!55!black
}
\setlist[itemize]{leftmargin=1.2em}
\setlist[enumerate]{leftmargin=1.4em}
\titleformat{\section}{\Large\bfseries}{\thesection}{0.6em}{}
\titleformat{\subsection}{\large\bfseries}{\thesubsection}{0.6em}{}
\titleformat{\subsubsection}{\normalsize\bfseries}{\thesubsubsection}{0.6em}{}
\newcolumntype{L}[1]{>{\raggedright\arraybackslash}p{#1}}
\setlength{\tabcolsep}{4pt}
\renewcommand{\arraystretch}{1.15}

\newcommand{\doc}[1]{\textbf{#1}}
\newcommand{\slide}[1]{\emph{#1}}
\newcommand{\term}[1]{\textbf{#1}}

\title{\textbf{Confronto teorico tra RAD, SDD, ODD e slide del corso}\\
\large Preparazione all'esame orale di Ingegneria del Software}
\author{AFAM Connect}
\date{A.A. 2025/2026}

\begin{document}
\maketitle
\tableofcontents
\newpage

\section{Obiettivo del documento}

Questo documento mette in relazione i tre elaborati di progetto \doc{RAD}, \doc{SDD} e
\doc{ODD} di AFAM Connect con gli argomenti teorici presenti nelle slide della professoressa.
Lo scopo non e' riscrivere i documenti, ma preparare risposte da orale: per ogni blocco
teorico viene spiegato dove compare nel progetto, perche' e' coerente con le slide e quale
frase conviene usare in sede di esame.

Fonti considerate:
\begin{itemize}
  \item \doc{Rad.pdf}: requisiti, casi d'uso, oggetti di analisi, diagrammi dinamici,
  diagrammi delle classi e interfacce.
  \item \doc{SDD.pdf}: architettura proposta, decomposizione in sottosistemi, repository,
  mapping hardware/software e persistenza.
  \item \doc{ODD\_AFAM\_Connect.pdf}: package, classi concrete, JavaFX/FXML, SQLite,
  JDBC, Singleton, DTO e diagrammi UML di dettaglio.
  \item Slide: raccolta requisiti, casi d'uso, requisiti non funzionali, classi e oggetti,
  analisi, sequence diagram, class diagram, system design, object design, design pattern e
  testing.
\end{itemize}

\section{Mappa sintetica: teoria e documenti}

{\small
\begin{longtable}{L{3.0cm}L{4.4cm}L{6.0cm}}
\toprule
\textbf{Tema teorico} & \textbf{Slide collegate} & \textbf{Dove compare nel progetto} \\
\midrule
\endhead
Raccolta requisiti & \slide{3. Raccolta e modellazione dei requisiti};
\slide{3.1 Descrivere i casi d'uso}; \slide{3.2 Requisiti non funzionali} &
Nel \doc{RAD}: obiettivo generale, sistema attuale/proposto, requisiti funzionali,
requisiti non funzionali, attori, casi d'uso e flussi. \\
Modello funzionale & Slide sui casi d'uso e sulla descrizione del flusso degli eventi &
Nel \doc{RAD}: vista degli attori, vista d'insieme dei casi d'uso, schede di dettaglio
per login, registrazione, gestione profilo, gestione contenuti, condivisione e
consultazione. \\
Modello di analisi & \slide{5. Analisi} &
Nel \doc{RAD}: oggetti \term{entity}, \term{boundary} e \term{control}; modelli dinamici;
modelli delle classi. \\
Sequence diagram & \slide{6. Diagrammi di Interazione} &
Nel \doc{RAD}: una sequenza per ciascun caso d'uso significativo, usata per collegare
attori, schermate, control e DBMS. \\
Class diagram & \slide{4. Classi e oggetti}; \slide{7. Diagrammi delle Classi} &
Nel \doc{RAD}: class diagram di analisi; nel \doc{ODD}: class diagram di dettaglio per
package entity, utility, interfacce e control. \\
System design & \slide{8. System Design} &
Nel \doc{SDD}: obiettivi del sistema, architettura proposta, decomposizione in
sottosistemi, mapping hardware/software, modello ER e modello relazionale. \\
Object design & \slide{8.1 Object Design} &
Nel \doc{ODD}: raffinamento verso classi Java, package, visibilita', interfacce, scelte
di librerie, persistenza e trade-off implementativi. \\
Repository architecture & \slide{8. System Design}, sezione sulle architetture repository &
Nel \doc{SDD} e nel \doc{ODD}: i sottosistemi non comunicano direttamente tra loro ma
passano dal \texttt{DatabaseManager}/DBMS. \\
Design pattern & \slide{10-Design\_patterns} &
Nel \doc{ODD}: uso esplicito di Singleton per \texttt{SceneManager},
\texttt{SessioneCorrente} e \texttt{DatabaseManager}; struttura assimilabile a
MVC/Boundary-Control-Entity e Repository. \\
Testing & \slide{10. Testing} &
Non e' un documento separato, ma il progetto fornisce basi testabili: requisiti, casi
d'uso, sequence diagram e classi control/entity da trasformare in test di unita',
integrazione, sistema e accettazione. \\
\bottomrule
\end{longtable}
}

\section{RAD e teoria dei requisiti}

\subsection{RAD come specifica dei requisiti}

Le slide sulla raccolta dei requisiti spiegano che la specifica dei requisiti descrive lo
scopo del sistema dal punto di vista di clienti, utenti e sviluppatori, e che svolge il ruolo
di contratto tra committente e team di sviluppo. Questo coincide direttamente con il
\doc{RAD} di AFAM Connect: il documento dichiara l'obiettivo di creare una piattaforma
per il portfolio multimediale certificato degli studenti AFAM, con gestione della privacy,
condivisioni temporizzate, codici di sblocco e protezione anti-download.

All'orale:
\begin{quote}
Il RAD rappresenta la vista utente del sistema. Non descrive ancora come il software sara'
implementato, ma cosa deve fare, con quali vincoli e in quali scenari d'uso.
\end{quote}

\subsection{Sistema attuale e sistema proposto}

Nelle slide la raccolta dei requisiti parte dal problema e dal sistema attuale: bisogna
capire come gli utenti svolgono oggi le attivita' e quale miglioramento introduce il nuovo
sistema. Il \doc{RAD} e lo \doc{SDD} fanno proprio questo:
\begin{itemize}
  \item il sistema attuale e' frammentato: documenti separati, file inviati via email,
  mancanza di repository centralizzata e poca tracciabilita';
  \item il sistema proposto centralizza profilo, contenuti artistici, privacy, link
  temporizzati, codici di sblocco e consultazione controllata.
\end{itemize}

Il collegamento teorico e' importante: il progetto non parte dalla tecnologia, ma da un
bisogno di dominio. La tecnologia appare dopo, nello \doc{SDD} e nell'\doc{ODD}.

\subsection{Requisiti funzionali}

Le slide distinguono i requisiti funzionali come servizi e comportamenti che il sistema
deve offrire. Nel \doc{RAD}, i requisiti funzionali sono organizzati per macro-attori:
\begin{itemize}
  \item \term{Studente}: autenticazione, registrazione, recupero password, gestione
  profilo, background artistico, categorie, upload file, modifica/eliminazione file,
  gestione privacy, generazione e revoca di link/codici.
  \item \term{Utente/Utente esterno}: ricerca studente, visualizzazione profilo,
  consultazione contenuti pubblici, accesso a contenuti privati con codice o link,
  fruizione protetta degli asset.
\end{itemize}

All'orale conviene sottolineare che un requisito funzionale in AFAM Connect non e' solo
"caricare file", ma include l'interazione completa: selezione categoria, tipologia,
memorizzazione, stato di privacy e successiva disponibilita' alla consultazione.

\subsection{Requisiti non funzionali}

Le slide sui requisiti non funzionali spiegano che essi non descrivono singoli servizi,
ma vincoli e qualita' complessive: usabilita', affidabilita', prestazioni, sicurezza,
privacy, standard, interoperabilita' e vincoli organizzativi. Il \doc{RAD} rispecchia questa
classificazione:
\begin{itemize}
  \item \term{Requisiti del prodotto}: usabilita', tempi di risposta, disponibilita',
  robustezza, protezione anti-download e gestione della privacy.
  \item \term{Requisiti organizzativi}: uso in contesto AFAM, gestione centralizzata,
  coerenza con flussi istituzionali.
  \item \term{Requisiti esterni}: privacy, sicurezza, possibilita' futura di integrazione
  con identita' digitali come SPID/eIDAS.
\end{itemize}

Punto da ricordare: nelle slide si dice che un requisito non funzionale puo' influenzare
l'architettura. Nel progetto accade davvero: privacy, robustezza e tracciabilita' portano
alla scelta di una repository centralizzata, alla gestione dei permessi nel DBMS, alla
temporizzazione dei link e al controllo tramite \texttt{DatabaseManager}.

\subsection{Attori, scenari e casi d'uso}

Le slide sulla raccolta dei requisiti indicano una sequenza tipica:
\begin{enumerate}
  \item identificare gli attori;
  \item identificare scenari;
  \item derivare casi d'uso;
  \item raffinare i casi d'uso;
  \item individuare relazioni;
  \item individuare requisiti non funzionali.
\end{enumerate}

Il \doc{RAD} segue questa logica. Gli attori principali sono lo Studente, l'Utente/Utente
esterno, il DBMS e il Tempo. I casi d'uso coprono autenticazione, profilo, contenuti,
condivisione tramite link, condivisione interna tramite codice, consultazione del profilo
e visualizzazione tramite link.

\subsection{Descrizione dei casi d'uso}

Le slide \slide{3.1 Descrivere i casi d'uso} chiariscono che il diagramma dei casi d'uso non
basta: ogni caso d'uso deve essere completato con nome, attori, precondizioni,
flusso base, flussi alternativi/eccezioni e postcondizioni. Il \doc{RAD} applica questa idea
nelle schede di dettaglio. Ad esempio:
\begin{itemize}
  \item \term{Login}: lo studente inserisce credenziali, il sistema interroga il DBMS,
  verifica, gestisce esito positivo o errore.
  \item \term{Gestisci Privacy Contenuti}: lo studente seleziona un file, sceglie
  Pubblico/Privato, il sistema aggiorna il DBMS e mostra conferma.
  \item \term{Caduta di connessione}: caso d'uso trasversale che descrive la gestione
  delle anomalie quando serve un'interazione transazionale con la base dati.
\end{itemize}

All'orale:
\begin{quote}
I casi d'uso del RAD non sono solo ovali nel diagramma. Sono specifiche testuali
verificabili: dicono quando il caso inizia, chi interagisce, quali passi avvengono e
quando il caso termina.
\end{quote}

\section{RAD e teoria dell'analisi}

\subsection{Il modello di analisi: funzionale, oggetti, dinamico}

Le slide di analisi spiegano che il modello di analisi e' formato da tre parti:
\begin{itemize}
  \item \term{modello funzionale}: casi d'uso e scenari;
  \item \term{modello degli oggetti}: classi, attributi, operazioni e relazioni visibili
  nel dominio;
  \item \term{modello dinamico}: sequence diagram e state machine, cioe' comportamento
  nel tempo.
\end{itemize}

Il \doc{RAD} contiene tutti e tre:
\begin{itemize}
  \item casi d'uso e attori nella sezione dei modelli di sistema;
  \item oggetti entity, boundary e control nella sezione del modello degli oggetti;
  \item diagrammi di sequenza per i casi d'uso principali;
  \item diagrammi delle classi per le funzionalita' individuate.
\end{itemize}

Questa e' una corrispondenza molto forte tra progetto e slide.

\subsection{Entity, boundary e control}

Le slide \slide{5. Analisi} distinguono tre tipi di oggetti:
\begin{itemize}
  \item \term{entity}: informazioni persistenti o concetti del dominio;
  \item \term{boundary}: interazioni tra attori e sistema, spesso schermate o form;
  \item \term{control}: coordinamento del caso d'uso, logica applicativa e passaggio tra
  boundary ed entity.
\end{itemize}

Nel \doc{RAD}:
\begin{itemize}
  \item \term{EntityStudente} rappresenta lo studente e i suoi dati principali;
  \item le \term{boundary} sono schermate come \texttt{SchermataLogin},
  \texttt{SchermataUpload}, \texttt{SchermataPrivacy}, \texttt{SchermataElencoLink};
  \item i \term{control} sono classi come \texttt{LoginControl},
  \texttt{UploadFileControl}, \texttt{GestioneFileControl},
  \texttt{GestisciLinkControl}, \texttt{CercaStudenteControl}.
\end{itemize}

Il punto teorico da dire bene e': il control non e' una schermata e non e' un dato del
dominio; e' l'oggetto che coordina il flusso del caso d'uso. Questa separazione riduce
l'accoppiamento e rende piu' chiara la responsabilita' di ciascuna classe.

\subsection{Dai casi d'uso agli oggetti}

Le slide mostrano che dagli scenari e dai casi d'uso si ricavano oggetti partecipanti.
Nel progetto questo passaggio si vede cosi':
\begin{itemize}
  \item dal caso d'uso \term{Upload File} nascono boundary di upload, control per
  gestire caricamento e dati/file persistiti;
  \item dal caso d'uso \term{Genera Link} nascono schermate di selezione contenuti,
  control di generazione link e dati persistenti del link;
  \item dal caso d'uso \term{Consulta Profilo Studente} nascono schermate di ricerca,
  control di consultazione e accesso ai dati dello studente e dei contenuti pubblici.
\end{itemize}

All'orale:
\begin{quote}
Abbiamo usato i casi d'uso come base per scoprire gli oggetti partecipanti. La vista
funzionale produce il flusso; il modello degli oggetti assegna responsabilita' e struttura.
\end{quote}

\section{Diagrammi UML: classi e sequenze}

\subsection{Sequence diagram}

Le slide sui diagrammi di interazione spiegano che i sequence diagram rappresentano
la sequenza temporale dei messaggi tra attori e oggetti. Sono usati in analisi per:
\begin{itemize}
  \item raffinare la descrizione dei casi d'uso;
  \item trovare oggetti partecipanti mancanti;
  \item distribuire il comportamento tra boundary, control ed entity;
  \item assegnare responsabilita' sotto forma di operazioni.
\end{itemize}

Nel \doc{RAD}, ogni sequence diagram rende concreto un caso d'uso. Per esempio, nel
login il flusso tipico e': Studente $\rightarrow$ schermata di login $\rightarrow$
\texttt{LoginControl} $\rightarrow$ DBMS $\rightarrow$ esito e navigazione. Nel caso
\term{Gestisci Privacy Contenuti}, il flusso mette in evidenza la schermata che raccoglie
la scelta e il control che aggiorna il DBMS.

Punto importante: nelle slide, nelle euristiche per disegnare i sequence diagram, viene
detto che la prima colonna dovrebbe essere l'attore che inizia il caso d'uso, la seconda
colonna una boundary e la terza un control object che gestisce il resto del caso d'uso.
Questo corrisponde a una logica di \term{sequence diagram a controllo centralizzato}:
dopo l'interazione iniziale dell'attore con la schermata, il flusso viene coordinato dal
control object.

Questa regola si ritrova nel progetto: l'attore non parla direttamente col DBMS, ma passa
dalla schermata e poi dalla logica di controllo. Ad esempio:
\begin{center}
\texttt{Studente} $\rightarrow$ \texttt{SchermataPrivacy} $\rightarrow$
\texttt{GestioneFileControl} $\rightarrow$ \texttt{DatabaseManager/DBMS}
\end{center}
In questo schema la boundary raccoglie l'input, mentre il control centralizza la gestione
del caso d'uso, decide quali operazioni invocare e coordina l'accesso ai dati. Non va
confuso con un unico controller globale dell'intero sistema: la centralizzazione e'
locale al singolo caso d'uso.

\subsection{Class diagram}

Le slide sui class diagram spiegano che il diagramma delle classi rappresenta classi,
attributi, operazioni, associazioni, molteplicita' e generalizzazioni. Nel progetto i
class diagram compaiono in due livelli:
\begin{itemize}
  \item nel \doc{RAD}: class diagram di analisi, ancora collegati ai casi d'uso e ai
  concetti del dominio;
  \item nell'\doc{ODD}: class diagram di dettaglio, collegati a package Java,
  interfacce grafiche, controller, utility e persistenza.
\end{itemize}

Esempio di collegamento teoria-progetto:
\begin{itemize}
  \item le slide dicono che gli attributi sono "piccole cose" interne alla classe, mentre le
  associazioni rappresentano relazioni con classi importanti;
  \item nel progetto, dati semplici come email, telefono, stato privacy o token sono
  trattabili come attributi/record, mentre studente, contenuto, link e codice sono concetti
  rilevanti della persistenza e vengono rappresentati nel modello dati.
\end{itemize}

\section{SDD e teoria del System Design}

\subsection{Che cosa aggiunge il System Design}

Le slide di System Design affermano che il modello di analisi descrive il sistema dal
punto di vista degli attori, ma non dice ancora come il sistema sara' organizzato
internamente. Il \doc{SDD} nasce proprio per colmare questo passaggio.

Nel \doc{SDD} di AFAM Connect compaiono:
\begin{itemize}
  \item architettura software attuale;
  \item architettura software proposta;
  \item requisiti minimi per l'utilizzo;
  \item decomposizione in sottosistemi;
  \item mappatura hardware/software;
  \item gestione dei dati persistenti;
  \item modello ER, modello relazionale e struttura delle tabelle.
\end{itemize}

All'orale:
\begin{quote}
Il RAD risponde a "cosa deve fare il sistema"; lo SDD risponde a "come lo
organizziamo a livello architetturale per soddisfare quei requisiti".
\end{quote}

\subsection{Design goals e requisiti non funzionali}

Le slide dicono che i design goals derivano dai requisiti non funzionali e guidano i
trade-off. In AFAM Connect i principali design goals sono:
\begin{itemize}
  \item \term{sicurezza e privacy}: accesso controllato a file pubblici/privati;
  \item \term{tracciabilita'}: link e codici hanno stato, scadenza e revoca;
  \item \term{semplicita' di distribuzione}: applicazione desktop locale;
  \item \term{robustezza}: gestione degli errori DBMS e riduzione della corruzione dati;
  \item \term{manutenibilita'}: sottosistemi separati per coesione funzionale.
\end{itemize}

Questi obiettivi spiegano le scelte successive: repository centralizzata, DBMS relazionale,
SQLite embedded, \texttt{DatabaseManager} centralizzato e package separati.

\subsection{Decomposizione in sottosistemi}

Le slide presentano la decomposizione in sottosistemi come prodotto centrale del System
Design. Lo \doc{SDD} suddivide AFAM Connect in cinque sottosistemi:
\begin{itemize}
  \item \term{Gestione Studente}: profilo, background, categorie, file, privacy;
  \item \term{Gestione Autenticazione}: login, registrazione, recupero password, logout,
  SPID/eIDAS come estensione possibile;
  \item \term{Gestione Condivisione}: link temporizzati, codici di sblocco, revoca;
  \item \term{Consultazione}: ricerca, profili, contenuti pubblici, contenuti con codice
  o link;
  \item \term{Gestione Operazioni Automatiche}: controllo della validita' temporale dei
  link.
\end{itemize}

Questa divisione segue la coesione funzionale: ogni sottosistema raggruppa oggetti e
operazioni che collaborano per una stessa area di responsabilita'.

\subsection{Architettura Repository}

Il legame piu' forte tra \doc{SDD}, \doc{ODD} e slide e' l'architettura Repository. Le slide
descrivono il Repository come uno stile in cui tutti i dati sono gestiti in un archivio
centrale accessibile ai componenti; i componenti non interagiscono direttamente tra loro,
ma tramite repository.

Nel progetto:
\begin{itemize}
  \item i sottosistemi applicativi non comunicano direttamente tra loro;
  \item il DBMS e il \texttt{DatabaseManager} sono il punto centrale di accesso ai dati;
  \item dati di studenti, contenuti, categorie, link, codici e associazioni vengono gestiti
  in modo persistente;
  \item il repository consente coerenza e separazione tra funzionalita'.
\end{itemize}

Vantaggi, secondo le slide, applicati al progetto:
\begin{itemize}
  \item sottosistemi piu' indipendenti;
  \item dati gestiti in modo consistente;
  \item minore necessita' di conoscenza reciproca tra moduli.
\end{itemize}

Svantaggi/rischi da citare:
\begin{itemize}
  \item il repository diventa un punto critico;
  \item errori o indisponibilita' del DBMS influenzano tutto il sistema;
  \item per questo il \doc{RAD} include il caso d'uso \term{Caduta di connessione} e
  l'\doc{ODD} prevede gestione di \texttt{SQLException} tramite messaggi di avviso.
\end{itemize}

\subsection{Mapping hardware/software}

Le slide indicano che lo SDD deve descrivere come i sottosistemi sono assegnati
all'hardware e alla piattaforma software. Nel progetto il sistema e' pensato come
applicazione desktop locale installata su un nodo AFAM, con:
\begin{itemize}
  \item interfaccia JavaFX/FXML;
  \item logica applicativa Java;
  \item SQLite embedded;
  \item storage locale dei contenuti, collegato ai record del database tramite percorso.
\end{itemize}

La scelta e' coerente con il requisito di semplicita' distributiva: SQLite non richiede un
server separato e riduce dipendenze di rete. Se la prof chiede il trade-off, la risposta e':
la soluzione attuale e' semplice e adatta a un'app desktop locale, ma, cosi' com'e',
e' meno adatta a scenari multiutente distribuiti o con alta concorrenza.

Questo limite riguarda pero' il \term{deployment attuale}, non l'idea di architettura
Repository in se'. La stessa logica Repository potrebbe essere mantenuta come architettura
interna di una soluzione piu' esterna di tipo client-server: i client, desktop o web,
invierebbero richieste a un server applicativo; il server conterrebbe i sottosistemi di
controllo e il livello repository; il DBMS sarebbe centralizzato e gestirebbe persistenza,
concorrenza e accesso multiutente. In questo scenario cambierebbe la mappatura
hardware/software, ma resterebbe valida l'idea centrale dello SDD: i sottosistemi non
comunicano direttamente tra loro, bensi' accedono ai dati tramite un punto di persistenza
comune.

All'orale si puo' dire:
\begin{quote}
Nel progetto attuale la repository e' locale, quindi la soluzione e' pensata per una
distribuzione desktop semplice. Se volessimo supportare pienamente la multiutenza,
potremmo mantenere il pattern Repository, ma collocarlo lato server dentro
un'architettura client-server, usando un DBMS centralizzato e politiche di concorrenza
adeguate.
\end{quote}

\subsection{Gestione della persistenza}

Le slide dicono che i dati persistenti sono quelli che sopravvivono alla singola esecuzione
del sistema e che la loro gestione influenza l'architettura. Lo \doc{SDD} dedica una sezione
alla persistenza:
\begin{itemize}
  \item modello ER;
  \item modello relazionale;
  \item tabelle \texttt{Studente}, \texttt{Background\_Artistico},
  \texttt{Categoria\_Artistica}, \texttt{Contenuto\_File},
  \texttt{Link\_Condivisione}, \texttt{Dettaglio\_Link},
  \texttt{Codice\_Sblocco}, \texttt{Dettaglio\_Codice}.
\end{itemize}

Il collegamento col \doc{RAD} e' diretto: i dati persistenti derivano dai requisiti e dai casi
d'uso. Se esiste un caso d'uso per generare link, deve esistere un dato persistente per
rappresentare link, scadenza, stato e contenuti associati.

\section{ODD e teoria dell'Object Design}

\subsection{Che cosa aggiunge l'Object Design}

Le slide di Object Design spiegano che questa fase raffina analisi e system design
aggiungendo dettagli implementativi: attributi mancanti, operazioni, tipi, visibilita',
eccezioni, librerie, mapping verso codice. L'\doc{ODD} di AFAM Connect fa esattamente
questo:
\begin{itemize}
  \item definisce package Java;
  \item descrive librerie e tecnologie concrete;
  \item documenta interfacce grafiche JavaFX/FXML;
  \item specifica SQLite e JDBC;
  \item introduce Singleton e DTO;
  \item collega i diagrammi UML alle classi effettive.
\end{itemize}

All'orale:
\begin{quote}
L'ODD traduce le scelte architetturali dello SDD in classi, package, librerie e
decisioni implementative. E' il ponte tra modello e codice.
\end{quote}

\subsection{Package e coerenza con i sottosistemi}

L'\doc{ODD} organizza il codice nel package principale \texttt{it.afam} e nei package:
\begin{itemize}
  \item \texttt{entity};
  \item \texttt{utility};
  \item \texttt{gestioneAutenticazione};
  \item \texttt{gestioneStudente};
  \item \texttt{gestioneCondivisione};
  \item \texttt{consultazione};
  \item \texttt{gestioneOperazioniAutomatiche}.
\end{itemize}

Questo e' coerente con lo \doc{SDD}: i package dell'ODD riflettono i sottosistemi del
System Design. La teoria dice che la decomposizione deve ridurre complessita' e
accoppiamento; qui ogni area funzionale ha le proprie boundary e i propri control.

\subsection{Boundary, control, entity nel codice}

Nel \doc{RAD} boundary/control/entity sono oggetti di analisi. Nell'\doc{ODD} diventano
struttura concreta:
\begin{itemize}
  \item le boundary sono realizzate dalla coppia formata da file FXML e classe controller
  JavaFX associata: il file FXML descrive il layout della schermata, mentre il controller
  JavaFX gestisce gli eventi dell'interfaccia, come click, input e navigazione;
  \item i control BCE veri e propri, come \texttt{LoginControl},
  \texttt{GestioneFileControl} o \texttt{UploadFileControl}, contengono invece la logica
  applicativa dei casi d'uso;
  \item \texttt{EntityStudente} rappresenta il dato principale caricato in sessione;
  \item \texttt{DatabaseManager} e DTO gestiscono l'accesso ai dati senza esporre
  direttamente tutta la persistenza alle interfacce.
\end{itemize}

Va quindi evitata una confusione terminologica: un \term{controller JavaFX} non coincide
con un \term{control object} del modello Boundary-Control-Entity. Il controller JavaFX fa
parte della boundary, perche' serve a gestire la schermata; il control object BCE coordina
invece il caso d'uso e dialoga con il livello dati.

Questa e' una buona risposta se viene chiesto come si passa dall'analisi al codice:
\begin{quote}
Abbiamo mantenuto la separazione boundary-control-entity anche nel design di dettaglio:
le boundary sono realizzate con FXML e controller JavaFX, i control BCE coordinano il
caso d'uso, il database conserva i dati persistenti.
\end{quote}

\subsection{JavaFX, FXML e SceneManager}

L'\doc{ODD} motiva l'uso di JavaFX e FXML per le interfacce. In termini teorici:
\begin{itemize}
  \item le schermate sono boundary object;
  \item FXML separa descrizione grafica e codice Java;
  \item \texttt{SceneManager} centralizza la navigazione tra schermate;
  \item la navigazione coerente mantiene allineati i flussi descritti nel \doc{RAD}.
\end{itemize}

Il \texttt{SceneManager} e' indicato come Singleton: mantiene un unico \texttt{Stage} e
cambia la \texttt{Scene} corrente. Questo evita che ogni schermata gestisca autonomamente
la finestra principale.

\subsection{SQLite, JDBC e DatabaseManager}

L'\doc{ODD} specifica SQLite come DBMS embedded e JDBC come tecnologia di accesso. Il
\texttt{DatabaseManager} centralizza le operazioni CRUD e funge da punto di accesso alla
repository. Questa scelta collega tre livelli:
\begin{itemize}
  \item \doc{RAD}: servono privacy, link, codici, contenuti e robustezza;
  \item \doc{SDD}: si sceglie una repository e si modellano dati persistenti;
  \item \doc{ODD}: si implementa la repository con SQLite, JDBC e \texttt{DatabaseManager}.
\end{itemize}

Risposta da orale:
\begin{quote}
SQLite e' coerente con l'applicazione desktop locale: semplifica distribuzione e
installazione. JDBC mantiene l'accesso al DBMS standardizzato dal lato Java.
\end{quote}

\subsection{Singleton e DTO}

Le slide sui design pattern danno un catalogo di soluzioni ricorrenti. Nel progetto si
riconoscono almeno:
\begin{itemize}
  \item \term{Singleton}: \texttt{SceneManager}, \texttt{SessioneCorrente},
  \texttt{DatabaseManager};
  \item \term{Repository}: architettura dei dati centralizzata;
  \item \term{MVC/BCE}: separazione tra interfaccia, controllo e dati;
  \item \term{DTO}: oggetti/record per trasferire solo i dati necessari tra DBMS, control
  e interfaccia.
\end{itemize}

Il Singleton serve quando deve esistere un solo punto di accesso a una risorsa globale:
la finestra principale, la sessione corrente o la connessione/punto di accesso al database.
Il rischio teorico da citare e' che un Singleton puo' aumentare dipendenza globale se usato
male; nel progetto e' giustificato per risorse naturalmente uniche.

\subsection{Trade-off dell'ODD}

Le slide di Object Design parlano di trade-off e raffinamento. Nel progetto i trade-off piu'
difendibili sono:
\begin{itemize}
  \item \term{SQLite embedded}: facile da distribuire, ma meno adatto a forte concorrenza.
  \item \term{Repository centralizzata}: coerenza e disaccoppiamento, ma punto critico
  in caso di errore DBMS.
  \item \term{JavaFX/FXML}: separazione UI/codice e sviluppo rapido, ma dipendenza dal
  framework JavaFX.
  \item \term{DTO}: meno esposizione della persistenza alle UI, ma piu' classi/record da
  mantenere.
\end{itemize}

\section{Dal documento al codice: cosa mostrare se la prof chiede evidenza pratica}

Una domanda possibile all'orale e': ``Questa cosa che avete scritto nel RAD, SDD o ODD
dove si vede nel codice?''. In quel caso bisogna rispondere facendo sempre lo stesso
ragionamento:
\begin{enumerate}
  \item prima nomino il concetto nel documento;
  \item poi indico il package o la classe in cui e' realizzato;
  \item infine spiego il ruolo di quel codice nel flusso.
\end{enumerate}

\subsection{Mappa rapida documento-codice}

{\footnotesize
\begin{longtable}{L{3.0cm}L{5.4cm}L{5.7cm}}
\toprule
\textbf{Cosa dice il documento} & \textbf{Dove mostrarlo nel codice} & \textbf{Come spiegarlo all'orale} \\
\midrule
\endhead
ODD: package coerenti con i sottosistemi &
\path{Codice/src/main/java/it/afam/gestioneAutenticazione},
\path{gestioneStudente}, \path{gestioneCondivisione},
\path{consultazione}, \path{gestioneOperazioniAutomatiche} &
I package Java riflettono la decomposizione in sottosistemi dello SDD: ogni area
funzionale ha interfacce e control dedicati. \\
RAD/Analisi: boundary-control-entity &
\path{interfaccia/SchermataLogin.java}, \path{SchermataLogin.fxml},
\path{control/LoginControl.java}, \path{entity/EntityStudente.java} &
La boundary e' realizzata da FXML piu' controller JavaFX; il control BCE e'
\texttt{LoginControl}; l'entity e' \texttt{EntityStudente}. \\
RAD: caso d'uso Login &
\path{SchermataLogin.fxml}, \path{SchermataLogin.java},
\path{LoginControl.java}, \path{DatabaseManager.java} &
Il flusso del caso d'uso si vede nel codice: click su Accedi, chiamata al control,
verifica credenziali sul repository, invio OTP e passaggio alla schermata di verifica. \\
SDD: architettura Repository &
\path{utility/DatabaseManager.java} e \path{utility/BoundaryDBMS.java} &
\texttt{DatabaseManager} e' il punto unico di accesso ai dati; \texttt{BoundaryDBMS}
incapsula la connessione JDBC/SQLite. \\
SDD: dati persistenti e modello relazionale &
\path{Codice/src/main/resources/schema.sql} &
Le tabelle dello SDD sono realizzate nello schema SQL: \texttt{Studente},
\texttt{Contenuto\_File}, \texttt{Link\_Condivisione}, \texttt{Codice\_Sblocco} e tabelle
di dettaglio. \\
ODD: Singleton &
\path{DatabaseManager.getIstanza()}, \path{SceneManager.getIstanza()},
\path{SessioneCorrente.getIstanza()} &
Queste classi hanno costruttore privato e metodo statico di accesso all'istanza:
gestiscono risorse uniche come repository, stage JavaFX e sessione corrente. \\
ODD: DTO &
\path{utility/dto/DatiStudente.java}, \path{DatiFile.java},
\path{DatiLink.java}, \path{DatiCodice.java} &
I DTO trasportano solo i dati necessari tra repository, control e boundary, senza
esporre direttamente tutta la struttura del database o l'entity completa. \\
RAD: caduta di connessione &
\path{ConnessioneException.java}, \path{BoundaryDBMS.java},
gestione \texttt{ConnessioneException} nelle schermate &
Il caso d'uso trasversale di errore DBMS viene tradotto in eccezione applicativa e in
messaggi di avviso all'utente, ad esempio nel login se il database non risponde. \\
RAD/SDD: scadenza temporale dei link &
\path{DatabaseManager.java}, \path{BoundaryTempo.java},
\path{ControlValiditaLink.java}, \path{Main.java} &
La scadenza a 14 giorni e' calcolata con \path{LocalDate.now().plusDays(14)};
poi un controllo automatico confronta la data corrente con la data di scadenza e marca
come scaduti i link non piu' validi. \\
RAD/SDD: SPID/eIDAS come estensione &
\path{SchermataLogin.java} e \path{LoginControl.java} &
Il codice mostra il pulsante e il messaggio di funzionalita' non disponibile: e' coerente
con l'idea di estensione futura, non con una funzionalita' gia' completata. \\
\bottomrule
\end{longtable}
}

\subsection{Struttura del progetto e posizione dei file}

Se la prof non si ferma ai documenti e chiede di vedere dove le scelte sono state
realizzate nel codice, conviene partire dalla struttura fisica del progetto. La cartella
principale del codice e':
\begin{center}
\path{C:/Users/Simone/Desktop/IDS/Codice}
\end{center}
Dentro questa cartella, la parte piu' importante per l'orale e':
\begin{itemize}
  \item \path{src/main/java/it/afam}: contiene le classi Java dell'applicazione;
  \item \path{src/main/resources/fxml}: contiene le schermate JavaFX in formato FXML;
  \item \path{src/main/resources/css}: contiene lo stile grafico dell'interfaccia;
  \item \path{src/main/resources/schema.sql}: contiene lo schema del database SQLite;
  \item \path{src/main/resources/email.properties}: contiene la configurazione tecnica
  usata per l'invio delle email.
\end{itemize}

Questa divisione e' utile da spiegare perche' separa il codice applicativo dalle risorse.
Le classi Java descrivono comportamento, logica e accesso ai dati; i file FXML descrivono
le schermate; lo schema SQL descrive la parte persistente del sistema. Quindi, quando
all'orale dici che il progetto segue una logica a livelli, puoi mostrarlo anche nella
struttura delle cartelle.

\subsubsection{Package Java principali}

Sotto \path{src/main/java/it/afam} ci sono i package che corrispondono alle aree
funzionali del sistema:
\begin{itemize}
  \item \path{gestioneAutenticazione}: contiene login, registrazione, recupero password,
  attivazione account e verifica dell'identita'. Qui si trovano classi come
  \texttt{LoginControl} e le schermate JavaFX collegate all'autenticazione.
  \item \path{gestioneStudente}: contiene profilo, dati personali, background artistico,
  gestione categorie, upload, privacy e gestione dei contenuti caricati dallo studente.
  \item \path{gestioneCondivisione}: contiene la logica dei link temporizzati, dei codici
  di sblocco e delle schermate per generare, visualizzare o revocare condivisioni.
  \item \path{consultazione}: contiene ricerca studenti, consultazione profili e
  visualizzazione dei contenuti pubblici o condivisi.
  \item \path{gestioneOperazioniAutomatiche}: contiene la logica automatica legata alla
  validita' dei link, cioe' operazioni non avviate direttamente da un click dell'utente.
  \item \path{entity}: contiene le entity principali, in particolare
  \texttt{EntityStudente}.
  \item \path{utility}: contiene classi trasversali usate da piu' sottosistemi, come
  \texttt{DatabaseManager}, \texttt{BoundaryDBMS}, \texttt{SceneManager},
  \texttt{SessioneCorrente}, \texttt{Utils} e i DTO.
\end{itemize}

La frase chiave e':
\begin{quote}
I package non sono casuali: riprendono la decomposizione in sottosistemi dello SDD e la
raffinano nell'ODD in classi concrete. Per questo posso passare dal documento al codice
seguendo i nomi delle responsabilita' funzionali.
\end{quote}

\subsubsection{Sottocartelle \texttt{interfaccia} e \texttt{control}}

Nei package funzionali ricorre spesso una divisione interna:
\begin{itemize}
  \item \path{interfaccia}: contiene le classi JavaFX che fanno da controller delle
  schermate FXML;
  \item \path{control}: contiene i control BCE, cioe' le classi che coordinano la logica
  dei casi d'uso.
\end{itemize}

Questa distinzione e' molto importante. Una classe dentro \path{interfaccia}, ad esempio
\texttt{SchermataLogin.java}, non e' il control BCE anche se in JavaFX viene chiamata
``controller'' della schermata. E' una boundary, perche' legge input dalla schermata,
mostra messaggi e reagisce agli eventi grafici. La logica applicativa vera, invece, sta in
classi come \texttt{LoginControl.java}, dentro \path{control}.

Quindi, se la prof chiede dove sta Boundary-Control-Entity nel codice, puoi rispondere
cosi':
\begin{quote}
Nel codice JavaFX la boundary e' composta dal file FXML piu' la classe JavaFX in
\texttt{interfaccia}; il control BCE sta invece nel package \texttt{control}; le entity e
i dati persistenti stanno in \texttt{entity}, nei DTO e nel livello repository/database.
\end{quote}

\subsubsection{Cartella \texttt{resources}: FXML, CSS e database}

La cartella \path{src/main/resources} raccoglie file che non sono classi Java ma sono
essenziali per l'esecuzione:
\begin{itemize}
  \item in \path{fxml} si trovano schermate come \texttt{SchermataLogin.fxml},
  \texttt{SchermataRegistrazione.fxml}, \texttt{SchermataUpload.fxml},
  \texttt{SchermataGenerazioneLink.fxml}, \texttt{SchermataRicerca.fxml} e molte altre;
  \item in \path{css} si trova \texttt{afam-theme.css}, cioe' la parte di presentazione
  grafica;
  \item \texttt{schema.sql} contiene la struttura delle tabelle del database;
  \item \texttt{email.properties} contiene parametri di configurazione per il servizio
  email.
\end{itemize}

I file FXML si collegano alle classi Java tramite l'attributo \texttt{fx:controller}. Ad
esempio, \texttt{SchermataLogin.fxml} dichiara quale classe JavaFX deve gestire gli eventi
della schermata. Questo e' il punto concreto in cui si vede l'unione tra risorsa grafica e
boundary Java.

\subsubsection{Percorso consigliato da mostrare alla prof}

Se durante l'orale devi mostrare rapidamente la struttura senza perderti, puoi seguire
questo ordine:
\begin{enumerate}
  \item apri \path{src/main/java/it/afam} e mostra i package principali, collegandoli ai
  sottosistemi dello SDD;
  \item apri un package semplice, ad esempio \path{gestioneAutenticazione}, e mostra la
  separazione tra \path{interfaccia} e \path{control};
  \item apri \path{src/main/resources/fxml/SchermataLogin.fxml} per mostrare la boundary
  grafica;
  \item apri \path{SchermataLogin.java} per mostrare il controller JavaFX della schermata;
  \item apri \path{LoginControl.java} per mostrare il control BCE del caso d'uso;
  \item apri \path{DatabaseManager.java} e \path{schema.sql} per mostrare repository e
  persistenza;
  \item apri \path{SessioneCorrente.java} per spiegare lo stato temporaneo della sessione.
\end{enumerate}

Questo percorso e' efficace perche' segue lo stesso ragionamento dei documenti: dal caso
d'uso del RAD si passa al sottosistema dello SDD, poi alle classi dell'ODD, e infine
all'implementazione concreta in Java, FXML e SQL.

\subsection{Come leggere questa mappa all'orale}

La tabella precedente non va studiata come un elenco di file da citare a memoria, ma come
una traccia per dimostrare che i documenti non sono scollegati dal codice. La risposta
piu' convincente e' sempre costruita su tre livelli:
\begin{itemize}
  \item \term{livello documentale}: nel RAD, SDD o ODD abbiamo dichiarato un requisito,
  una scelta architetturale o una scelta di object design;
  \item \term{livello strutturale}: nel codice questa scelta si ritrova nei package, nelle
  classi, negli FXML, nello schema SQL o nelle utility;
  \item \term{livello comportamentale}: durante l'esecuzione, quelle classi collaborano
  per realizzare un caso d'uso concreto.
\end{itemize}

Per esempio, se la prof chiede dove si vede il modello Boundary-Control-Entity, non basta
dire ``nel package interfaccia e nel package control''. Bisogna spiegare il percorso:
\begin{quote}
Nel RAD abbiamo modellato i casi d'uso con boundary, control ed entity. Nel codice questa
separazione si vede perche' la schermata FXML e il suo controller JavaFX rappresentano la
boundary, la classe \texttt{LoginControl} rappresenta il control BCE del caso d'uso Login,
mentre \texttt{EntityStudente} e il livello di persistenza rappresentano i dati gestiti dal
sistema.
\end{quote}

Quindi il codice va mostrato come una realizzazione progressiva dei documenti:
\begin{center}
\textbf{RAD} \(\rightarrow\) \textbf{SDD} \(\rightarrow\) \textbf{ODD}
\(\rightarrow\) \textbf{codice}
\end{center}
Il RAD identifica cosa deve accadere, lo SDD decide l'organizzazione architetturale,
l'ODD raffina in classi e tecnologie, il codice implementa quelle decisioni.

\subsection{Spiegazione dettagliata dei collegamenti principali}

\subsubsection{Boundary realizzate con FXML e controller JavaFX}

Nel RAD, una boundary e' un oggetto concettuale che rappresenta il punto di interazione
tra attore e sistema. Nel codice JavaFX, pero', questa boundary viene realizzata da due
elementi:
\begin{itemize}
  \item il file \texttt{.fxml}, che descrive il layout grafico;
  \item la classe JavaFX associata tramite \texttt{fx:controller}, che gestisce gli eventi
  della schermata.
\end{itemize}

Per esempio, \texttt{SchermataLogin.fxml} contiene campi, pulsanti e collegamento al
controller; \texttt{SchermataLogin.java} contiene il metodo \texttt{cliccaAccedi()}, cioe'
la reazione al click dell'utente. Questa classe appartiene ancora alla boundary:
\texttt{LoginControl.java} e' invece il control BCE, perche' coordina la logica del caso
d'uso.

Frase pronta:
\begin{quote}
Nel codice JavaFX la boundary non e' solo il file FXML: e' la coppia FXML piu' controller
JavaFX. Il controller JavaFX gestisce eventi dell'interfaccia, mentre il control BCE
gestisce la logica applicativa del caso d'uso.
\end{quote}

\subsubsection{Control BCE e logica applicativa}

I control BCE sono le classi nei package \texttt{control}, ad esempio
\texttt{LoginControl}, \texttt{UploadFileControl}, \texttt{GestioneFileControl} e
\texttt{GestisciLinkControl}. Il loro compito non e' disegnare schermate, ma coordinare
il caso d'uso. Nel Login, \texttt{LoginControl}:
\begin{itemize}
  \item riceve email e password dalla boundary;
  \item usa \texttt{Utils.hashPassword()} per confrontare la password hashata;
  \item chiama \texttt{DatabaseManager.verificaCredenziali()};
  \item genera e invia il codice OTP;
  \item dopo la verifica del codice, aggiorna \texttt{SessioneCorrente}.
\end{itemize}

Questo e' un esempio concreto di sequence a controllo centralizzato: la boundary avvia il
flusso, ma il control coordina il resto.

\subsubsection{Repository e DatabaseManager}

Lo \doc{SDD} descrive una architettura Repository. Nel codice questa scelta si vede in
\texttt{DatabaseManager.java}, che funge da punto unico di accesso ai dati. I control non
scrivono direttamente query JDBC sparse nel codice dell'interfaccia: chiamano metodi del
\texttt{DatabaseManager}. Questo rende piu' chiara la separazione:
\begin{itemize}
  \item i control decidono cosa fare nel caso d'uso;
  \item \texttt{DatabaseManager} sa come leggere e aggiornare i dati;
  \item \texttt{BoundaryDBMS} incapsula il dettaglio tecnico della connessione JDBC a
  SQLite.
\end{itemize}

Il collegamento con lo SDD e' quindi forte: i sottosistemi non comunicano tra loro in modo
diretto, ma condividono lo stato persistente attraverso il repository. In una futura
architettura client-server, questa stessa idea potrebbe rimanere valida spostando il
repository lato server.

\subsubsection{Schema SQL e modello relazionale}

La parte dello SDD relativa al modello relazionale si vede in \texttt{schema.sql}. Qui non
ci sono solo tabelle generiche, ma le tabelle che corrispondono ai concetti del progetto:
\texttt{Studente}, \texttt{Background\_Artistico}, \texttt{Categoria\_Artistica},
\texttt{Contenuto\_File}, \texttt{Link\_Condivisione}, \texttt{Codice\_Sblocco} e le tabelle
di dettaglio per collegare link/codici ai file.

Questo e' utile se la prof chiede se il modello ER/relazionale dello SDD e' stato davvero
implementato. La risposta e':
\begin{quote}
Si', lo schema relazionale e' realizzato nel file \texttt{schema.sql}; il
\texttt{DatabaseManager} inizializza lo schema al primo avvio e poi usa quelle tabelle
attraverso metodi DAO.
\end{quote}

\subsubsection{Utilizzo concreto di EntityStudente}

\texttt{EntityStudente} va introdotta nel documento perche' e' l'entity principale del
progetto: rappresenta lo studente come oggetto del dominio applicativo. Nel codice si trova
in:
\begin{center}
\path{src/main/java/it/afam/entity/EntityStudente.java}
\end{center}
e contiene attributi coerenti con la tabella \texttt{Studente} dello \doc{SDD}, come
\texttt{id}, \texttt{nome}, \texttt{cognome}, \texttt{codiceFiscale}, \texttt{email},
\texttt{password}, \texttt{telefono}, \texttt{percorsoFoto} e la lista dei dati di
background artistico.

Il punto importante da dire all'orale e' che \texttt{EntityStudente} non e' semplicemente
una tabella e non e' nemmeno un local storage. E' la rappresentazione a oggetti dello
studente quando il sistema deve lavorare sui suoi dati in memoria. La fonte persistente
rimane il database SQLite, mentre \texttt{EntityStudente} viene creata, riempita o letta
dal \texttt{DatabaseManager}.

Nel codice l'utilizzo si vede soprattutto in questi punti:
\begin{itemize}
  \item in \texttt{RegistrazioneControl.java}, durante la registrazione, viene creato uno
  \texttt{studenteProvvisorio} di tipo \texttt{EntityStudente}; dopo la verifica, viene
  salvato tramite \texttt{DatabaseManager.salvaStudente()};
  \item in \texttt{LoginControl.java}, il repository restituisce un
  \texttt{EntityStudente} quando le credenziali sono corrette; dopo l'OTP, pero', in
  \texttt{SessioneCorrente} viene salvato soprattutto l'\texttt{id} dello studente, non
  l'intera entity;
  \item in \texttt{GestioneProfiloControl.java}, \texttt{EntityStudente} viene usata per
  caricare e aggiornare i dati personali dello studente autenticato;
  \item in \texttt{VisualizzaProfiloControl.java} e nelle schermate di consultazione, puo'
  essere usata per mostrare il profilo di uno studente cercato;
  \item in \texttt{DatabaseManager.java}, il metodo interno di mapping costruisce un
  \texttt{EntityStudente} a partire dal \texttt{ResultSet} del database.
\end{itemize}

Questa distinzione e' utile per evitare una risposta imprecisa. Se la prof chiede se
\texttt{EntityStudente} funge da local storage, la risposta migliore e':
\begin{quote}
No, non proprio. \texttt{EntityStudente} e' l'entity del dominio che rappresenta i dati
dello studente quando vengono caricati dal database o preparati per il salvataggio. Lo
stato temporaneo della sessione e' invece in \texttt{SessioneCorrente}, mentre la
persistenza vera resta nel database gestito da \texttt{DatabaseManager}.
\end{quote}

Il ruolo di \texttt{EntityStudente} e' quindi quello di fare da ponte tra il modello dati
persistente e la logica applicativa: il database conserva i record, il repository li mappa
in oggetti, i control li usano per realizzare i casi d'uso, e le boundary mostrano solo le
informazioni necessarie all'utente.

\texttt{SessioneCorrente} va tenuta distinta: non modella lo studente completo, ma conserva
solo lo stato temporaneo utile durante la navigazione, come id dello studente autenticato,
categoria corrente, file corrente e profilo consultato. Anche per questo non va presentata
come local storage persistente: viene azzerata al logout e si perde alla chiusura
dell'applicazione.

\subsubsection{Scadenza temporale dei link}

Un punto molto utile da mostrare nel codice riguarda la scadenza dei link dopo 14 giorni.
Qui si collegano bene RAD, SDD, ODD e implementazione:
\begin{itemize}
  \item nel \doc{RAD} il sistema prevede link temporizzati e quindi un vincolo funzionale
  sulla validita' nel tempo;
  \item nello \doc{SDD} questa scelta richiede dati persistenti, perche' il sistema deve
  memorizzare data di creazione, data di scadenza e stato del link;
  \item nell'\doc{ODD} la responsabilita' viene tradotta in classi concrete:
  \texttt{DatabaseManager}, \texttt{BoundaryTempo} e \texttt{ControlValiditaLink};
  \item nel codice la scadenza viene implementata usando le API \texttt{java.time}, in
  particolare \texttt{LocalDate}.
\end{itemize}

Nel file \texttt{DatabaseManager.java}, il metodo \texttt{salvaLink()} inserisce il link
nella tabella \texttt{Link\_Condivisione}. In quel momento salva:
\begin{itemize}
  \item \texttt{data\_creazione} come data corrente;
  \item \texttt{data\_scadenza} come data corrente piu' 14 giorni;
  \item \texttt{stato\_link} inizialmente uguale ad \texttt{Attivo}.
\end{itemize}

La parte chiave e':
\begin{center}
\path{LocalDate.now().plusDays(14)}
\end{center}
Quindi non e' una scadenza scritta a mano: viene calcolata dal sistema a partire dalla
data corrente.

Non e' stata introdotta una classe \texttt{Time} custom perche' non era necessaria:
Java fornisce gia' le API \texttt{java.time}, che sono piu' adatte e piu' sicure rispetto
a una gestione manuale delle date. Nel progetto, quindi, la responsabilita' del calcolo
temporale e' affidata a \texttt{LocalDate}, mentre \texttt{BoundaryTempo} serve solo a
isolare l'accesso alla data corrente e a rappresentare in modo chiaro l'interazione con
l'attore concettuale Tempo. Questa scelta e' coerente con il design: si evita di creare
una classe inutile e si mantiene comunque separato il punto in cui il sistema legge
l'orologio.

Poi entra in gioco il package \texttt{gestioneOperazioniAutomatiche}. La classe
\texttt{BoundaryTempo} rappresenta l'interfaccia verso l'attore concettuale Tempo, cioe'
l'orologio del sistema, e restituisce la data corrente con \texttt{LocalDate.now()}.
\texttt{ControlValiditaLink} recupera dal repository i link ancora attivi, legge la loro
data di scadenza e confronta:
\begin{center}
\texttt{oggi.isAfter(scadenza)}
\end{center}
Se la data corrente e' successiva alla scadenza, il control chiama
\texttt{DatabaseManager.disattivaLink()} e il link viene marcato come \texttt{Scaduto}.

Infine, in \texttt{Main.java}, il controllo viene avviato all'apertura dell'applicazione e
poi pianificato ogni giorno tramite \texttt{ScheduledExecutorService}. Questo significa che
la validita' del link non viene controllata solo quando lo studente guarda la lista dei
link, ma anche tramite un'operazione automatica periodica.

Frase pronta:
\begin{quote}
La scadenza dei link e' implementata con le API \texttt{java.time}. Quando genero un link,
\texttt{DatabaseManager} salva la data corrente e la data di scadenza calcolata con
\texttt{LocalDate.now().plusDays(14)}. Non abbiamo creato una classe \texttt{Time} custom
perche' Java offre gia' \texttt{java.time}; \texttt{BoundaryTempo} serve solo a isolare la
lettura della data corrente. Poi \texttt{ControlValiditaLink} confronta la data corrente
con la scadenza e marca il link come \texttt{Scaduto} se i 14 giorni sono stati superati.
\end{quote}

\subsubsection{DTO}

I DTO, come \texttt{DatiStudente}, \texttt{DatiFile}, \texttt{DatiLink} e
\texttt{DatiCodice}, servono a trasferire solo le informazioni necessarie tra database,
control e boundary. Sono utili perche' evitano di esporre sempre l'intera entity o tutta la
struttura della tabella. Per esempio, in una schermata di ricerca puo' bastare
\texttt{DatiStudente} con id, nome e cognome, senza caricare password, telefono o altri
dati non necessari.

Frase pronta:
\begin{quote}
I DTO sono oggetti di trasferimento dati: rendono piu' pulito il passaggio di informazioni
tra repository, control e boundary, mantenendo separati database, logica applicativa e
interfaccia.
\end{quote}

\subsection{Esempio guidato: dal caso d'uso Login al codice}

Se la prof chiede di mostrare un caso d'uso nel codice, il Login e' uno degli esempi piu'
chiari.

\begin{enumerate}
  \item Nel \doc{RAD}, il Login e' un caso d'uso della Gestione Autenticazione.
  \item Nel file \texttt{SchermataLogin.fxml} si vede la boundary grafica: campi email e
  password, pulsante Accedi e collegamento \texttt{fx:controller} alla classe JavaFX.
  \item In \texttt{SchermataLogin.java}, il metodo \texttt{cliccaAccedi()} gestisce l'evento
  dell'interfaccia: prende email/password, avvia il task e invoca il control.
  \item In \texttt{LoginControl.java}, il metodo \texttt{accedi()} contiene la logica del
  caso d'uso: hash della password, verifica credenziali, generazione OTP e invio email.
  \item In \texttt{DatabaseManager.java}, il metodo \texttt{verificaCredenziali()} interroga
  la tabella \texttt{Studente}.
  \item In \texttt{SessioneCorrente.java}, dopo verifica OTP, viene salvato l'id dello
  studente autenticato.
\end{enumerate}

Risposta pronta:
\begin{quote}
Il caso d'uso Login del RAD si ritrova nel codice attraversando boundary, control e
repository: la schermata FXML raccoglie i dati, il controller JavaFX gestisce il click,
\texttt{LoginControl} coordina la logica applicativa, \texttt{DatabaseManager} verifica le
credenziali sul DBMS e \texttt{SessioneCorrente} mantiene lo stato dello studente
autenticato.
\end{quote}

\subsection{Altri esempi rapidi da mostrare}

Dopo il Login, gli altri collegamenti da mostrare senza ripetere tutta la teoria sono:
\begin{itemize}
  \item \term{Repository}: partire da \texttt{DatabaseManager}; la catena da spiegare e'
  \texttt{Control BCE} $\rightarrow$ \texttt{DatabaseManager} $\rightarrow$
  \texttt{BoundaryDBMS} $\rightarrow$ \texttt{SQLite/schema.sql}.
  \item \term{Modello dati}: aprire \texttt{schema.sql}; le tabelle piu' importanti sono
  \texttt{Studente}, \texttt{Contenuto\_File}, \texttt{Link\_Condivisione},
  \texttt{Codice\_Sblocco} e le relative tabelle di dettaglio.
  \item \term{Singleton}: citare \texttt{SceneManager}, \texttt{DatabaseManager} e
  \texttt{SessioneCorrente}; sono risorse uniche per navigazione, repository e stato
  temporaneo della sessione.
\end{itemize}

Questi esempi vanno usati come conferma pratica di concetti gia' spiegati: repository,
persistenza e pattern non sono solo nel documento, ma hanno classi e file precisi nel
codice.

\subsection{Domande pronte sul codice}

\subsubsection*{Dove si vede la separazione Boundary-Control-Entity?}

Nel Login: \texttt{SchermataLogin.fxml} e \texttt{SchermataLogin.java} sono boundary,
\texttt{LoginControl.java} e' il control BCE, \texttt{EntityStudente.java} rappresenta
l'entity, mentre \texttt{DatabaseManager.java} gestisce la persistenza.

\subsubsection*{\texttt{EntityStudente} e' una specie di local storage?}

No. \texttt{EntityStudente} e' il modello a oggetti dello studente quando i dati vengono
caricati dal database o preparati per il salvataggio. La persistenza vera resta in SQLite;
lo stato temporaneo della navigazione e' in \texttt{SessioneCorrente}. Quindi non e' un
local storage, ma una entity del dominio applicativo.

\subsubsection*{Dove si vede il Repository?}

In \texttt{DatabaseManager.java}. I control non accedono direttamente a JDBC o SQLite,
ma usano i metodi del \texttt{DatabaseManager}. La connessione tecnica e' incapsulata in
\texttt{BoundaryDBMS.java}.

\subsubsection*{Dove si vede la mappatura dello SDD in codice?}

Nei package sotto \texttt{it.afam}: i nomi dei package corrispondono ai sottosistemi dello
SDD, mentre \texttt{schema.sql} corrisponde alla parte di gestione dei dati persistenti.

\subsubsection*{Dove si vede che SPID/eIDAS non e' completo?}

In \texttt{SchermataLogin.java}, il metodo associato al pulsante SPID/eIDAS mostra un
messaggio di funzionalita' non disponibile. In \texttt{LoginControl.java} il commento
specifica che il login federato e' un'estensione futura. Quindi va detto chiaramente:
il progetto lo prevede come estensione, ma non lo implementa completamente.

\subsubsection*{Dove si vede la gestione della caduta di connessione?}

In \texttt{BoundaryDBMS.java}, gli errori JDBC vengono convertiti in
\texttt{ConnessioneException}; nelle schermate, ad esempio \texttt{SchermataLogin.java},
l'eccezione viene intercettata e trasformata in messaggio di avviso per l'utente.

\subsubsection*{Dove si vede la scadenza dei link dopo 14 giorni?}

In \texttt{DatabaseManager.java}, nel metodo \texttt{salvaLink()}, la data di scadenza
viene calcolata con \texttt{LocalDate.now().plusDays(14)} e salvata nella tabella
\texttt{Link\_Condivisione}. \texttt{BoundaryTempo} fornisce la data corrente e
\texttt{ControlValiditaLink} confronta oggi con la scadenza: se i 14 giorni sono superati,
chiama \texttt{DatabaseManager.disattivaLink()}. In \texttt{Main.java} il controllo viene
eseguito all'avvio e poi pianificato ogni giorno.

\section{Testing: come collegare il progetto alle slide}

Il testing non e' uno dei tre documenti principali, ma le slide sono molto collegabili al
progetto. Le slide dicono che:
\begin{itemize}
  \item i test dimostrano la presenza di difetti, non la loro assenza;
  \item per software personalizzato deve esserci almeno una prova per ogni requisito;
  \item i casi d'uso possono essere base per i test di sistema;
  \item i sequence diagram documentano componenti e interazioni da testare;
  \item i test automatici hanno setup, chiamata e asserzione.
\end{itemize}

\subsection{Test basati su requisiti}

Dal \doc{RAD} si possono derivare test per ogni requisito:
\begin{itemize}
  \item login con credenziali corrette e errate;
  \item registrazione con dati validi/non validi;
  \item upload di file ammessi/non ammessi;
  \item modifica privacy da Pubblico a Privato;
  \item generazione link con scadenza;
  \item revoca link/codice;
  \item accesso a contenuto privato con codice valido/non valido;
  \item tentativo di fruizione con link scaduto.
\end{itemize}

\subsection{Test basati su casi d'uso e sequence diagram}

Le slide dicono che i use case testing sono utili per forzare l'interazione tra componenti.
Nel progetto, un test di sistema su \term{Genera Link} dovrebbe verificare:
\begin{enumerate}
  \item lo studente seleziona contenuti;
  \item il control genera token e scadenza;
  \item il DBMS salva link e dettagli;
  \item la schermata mostra conferma;
  \item il link consente accesso ai contenuti selezionati fino alla scadenza.
\end{enumerate}

Questo deriva direttamente dal sequence diagram e dal flusso del caso d'uso.

\subsection{Test di unita', componenti, sistema e accettazione}

Applicazione al progetto:
\begin{itemize}
  \item \term{Test di unita'}: metodi di \texttt{Utils}, generazione token/codici, validazione
  email/password, hashing, metodi DAO isolati.
  \item \term{Test di componenti}: interazione tra control e \texttt{DatabaseManager}, ad
  esempio upload o modifica privacy.
  \item \term{Test di sistema}: flussi completi come registrazione, login, upload,
  condivisione e consultazione.
  \item \term{Test di accettazione}: scenari reali con studente AFAM e valutatore esterno,
  per decidere se il sistema soddisfa il bisogno operativo.
\end{itemize}

\section{Domande orali probabili e risposte pronte}

\subsection*{1. Che differenza c'e' tra RAD, SDD e ODD?}

Il \doc{RAD} descrive il problema e i requisiti dal punto di vista degli utenti: attori,
casi d'uso, requisiti funzionali/non funzionali, oggetti di analisi e modelli dinamici.
Lo \doc{SDD} descrive la soluzione architetturale: sottosistemi, repository, mapping
hardware/software e persistenza. L'\doc{ODD} raffina la soluzione in classi, package,
librerie, pattern e scelte implementative.

\subsection*{2. Perche' il RAD e' collegato alle slide sui requisiti?}

Perche' segue la struttura teorica della specifica dei requisiti: scopo del sistema, sistema
attuale, sistema proposto, requisiti funzionali, requisiti non funzionali, attori, casi d'uso
e modelli di sistema. Inoltre usa il linguaggio dell'utente e non descrive ancora i dettagli
interni dell'implementazione.

\subsection*{3. Perche' i requisiti non funzionali sono importanti nel progetto?}

Perche' guidano le scelte architetturali. Privacy, robustezza, sicurezza, tempi di risposta
e protezione anti-download non sono semplici funzioni: influenzano repository, DBMS,
gestione dei link temporizzati, codici di sblocco, visualizzatore protetto e gestione degli
errori.

\subsection*{4. Che ruolo hanno entity, boundary e control?}

Sono una classificazione di analisi. Le entity rappresentano dati e concetti persistenti,
le boundary rappresentano l'interazione con l'utente, i control coordinano i casi d'uso.
Nel progetto questa separazione appare prima nel \doc{RAD} e poi viene mantenuta
nell'\doc{ODD} tramite package di interfaccia e package di controllo.

\subsection*{5. Perche' avete scelto una repository?}

Perche' il sistema ruota attorno a dati condivisi e persistenti: studenti, contenuti,
privacy, link, codici e categorie. La repository consente ai sottosistemi di non parlarsi
direttamente e di accedere ai dati tramite un punto centralizzato. Il vantaggio e'
disaccoppiamento e coerenza; lo svantaggio e' che il DBMS diventa un punto critico,
percio' e' stata prevista la gestione della caduta di connessione.

\subsection*{6. Come si passa dai casi d'uso ai sequence diagram?}

Il caso d'uso descrive il flusso testuale. Il sequence diagram lo rende operativo nel tempo:
mostra attore, boundary, control, DBMS e messaggi. In questo modo si scoprono oggetti
partecipanti, responsabilita' e operazioni necessarie.

\subsection*{7. Come si passa dallo SDD all'ODD?}

Lo \doc{SDD} dice quali sottosistemi esistono e come sono collegati; l'\doc{ODD} li traduce
in package e classi concrete. Ad esempio, il sottosistema Gestione Studente diventa un
package con interfacce FXML e control Java; la repository dello SDD diventa
\texttt{DatabaseManager}, SQLite, JDBC e DTO.

\subsection*{8. Dove compaiono i design pattern?}

Il progetto usa Singleton per \texttt{SceneManager}, \texttt{SessioneCorrente} e
\texttt{DatabaseManager}. Inoltre l'architettura complessiva richiama Repository e la
separazione Boundary-Control-Entity/MVC. Questi pattern aiutano a comunicare scelte
ricorrenti e a giustificare la struttura.

\subsection*{9. Come testereste il sistema usando le slide sul testing?}

Partirei dai requisiti e dai casi d'uso del \doc{RAD}. Ogni requisito deve avere almeno un
test. I sequence diagram aiutano a costruire test di sistema per le interazioni tra
schermate, control e database. I metodi di utilita' e DAO sono adatti a test di unita';
flussi come login, upload, privacy e generazione link sono adatti a test di integrazione
e sistema; gli scenari con studente e valutatore sono test di accettazione.

\section{Schema finale da memorizzare}

\begin{center}
\begin{tabular}{p{3cm}p{9.8cm}}
\toprule
\textbf{Documento} & \textbf{Frase chiave da orale} \\
\midrule
RAD & Specifica cosa deve fare il sistema dal punto di vista degli utenti: requisiti,
casi d'uso, attori e modelli di analisi. \\
SDD & Decide l'organizzazione architetturale: sottosistemi, repository, mapping
hardware/software e dati persistenti. \\
ODD & Raffina l'architettura in classi, package, pattern, tecnologie e scelte
implementative. \\
\bottomrule
\end{tabular}
\end{center}

\section{Conclusione}

Il progetto AFAM Connect segue in modo abbastanza lineare il percorso presentato nelle
slide: dalla raccolta dei requisiti si passa all'analisi con casi d'uso, oggetti e sequence
diagram; dall'analisi si passa al system design con sottosistemi, repository e persistenza;
dal system design si passa all'object design con package Java, JavaFX, SQLite, JDBC,
Singleton e DTO. Per l'orale conviene quindi non studiare RAD, SDD e ODD come tre
documenti separati, ma come tre livelli successivi dello stesso ragionamento progettuale.

\end{document}
